% ============================================================
% Cheat sheet, argomento: il pattern Composite (dispensa, capitolo 15)
% ============================================================
\argomento{Composite}{Dispensa: capitolo 15. Pattern strutturale della GoF}[Pattern Composite]

\scheda{Carta d'identità}{
  \pintento{Comporre gli oggetti in \textbf{strutture ad albero} (gerarchie tutto-parte) e permettere al client di
    trattare \textbf{allo stesso modo} oggetti singoli e oggetti composti.}
  \ptesto{«gerarchia», «albero», «un elemento può contenere altri elementi, che a loro volta\ldots», «cartelle e
    file», «menu e sottomenu», «gruppi e sottogruppi», «pannelli che contengono pulsanti e altri pannelli»,
    «un'operazione sul tutto si applica a tutte le parti», «trattare in modo uniforme».}
  \pcodice{un'interfaccia (o classe astratta) Component; una classe \textbf{foglia} che la implementa; una classe
    \textbf{composta} che la implementa \textbf{e} ha una \textbf{collezione dello stesso tipo}
    (\code{List<AbstractFile>}), con \code{add} e \code{remove}; il metodo del composto chiama lo \textbf{stesso
    metodo} su tutti i figli (ricorsione), e il client lo chiama una volta sola, sulla radice.}
  \pdiagramma{\code{Leaf} e \code{Composite} realizzano (o estendono) \code{Component}; da \code{Composite} parte
    un'\textbf{aggregazione con molteplicità \code{*}} che torna verso \code{Component}: un composto contiene
    componenti, quindi anche altri composti.}
  \pvarianti{\textbf{trasparenza}: \code{add}, \code{remove}, \code{getChild} dichiarati nel Component (il client tratta
    tutti allo stesso modo, ma deve gestire \code{add} su una foglia, per esempio con un'eccezione) oppure
    \textbf{sicurezza}: solo nel Composite (il compilatore impedisce di aggiungere figli a una foglia). Nelle GUI di
    Java: un \code{Container} è un \code{Component} che contiene \code{Component}.}
  \pconfondere{il \textbf{Decorator} (\vv{decorator-struttura}): anche lui implementa il Component e ne contiene uno,
    ma \textbf{uno solo}, e aggiunge comportamento; il Composite ne contiene \textbf{molti} e combina i risultati. Una
    semplice associazione riflessiva (\code{List<Persona> figli}): senza un'interfaccia comune fra foglie e composti
    non è un Composite.}
}

\begin{struttura}{composite-struttura}{La struttura}
\begin{tikzpicture}
  \node[classe1=1.3cm] (c) at (-2.8,0) {Client};
  \umlclass[3.4cm]{co}{(1.0,0)}{\textit{Component}}{}{+ operation()\\+ add(c: Component)\\+ remove(c: Component)\\+ getChild(i: int)}
  \umlclass[2.0cm]{l}{(-1.2,-3.0)}{Leaf}{}{+ operation()}
  \umlclass[3.4cm]{cm}{(2.6,-3.0)}{Composite}{}{+ operation()\\+ add(c: Component)\\+ remove(c: Component)\\+ getChild(i: int)}
  \draw[navig] (c.east) -- (co.west);
  \draw[gen] (l.north) |- ++(0,0.55) -| ([xshift=-0.6cm]co.south);
  \draw[gen] (cm.north) |- ++(0,0.4) -| ([xshift=0.6cm]co.south);
  \draw[aggr] (cm.east) -- ++(0.5,0) |- (co.east);
  \node[molt, anchor=south west] at ($(co.east)+(0.1,0.05)$) {*};
  \node[ruolo, anchor=north west] at ($(co.east)+(0.1,-0.05)$) {figli};
  \node[notauml, anchor=north, text width=2.6cm] (n) at (2.6,-4.6) {for all g in figli:\\\hspace*{1em}g.operation()};
  \draw[tratt] (cm.south) -- (n.north);
\end{tikzpicture}
\tcblower
\textbf{Component}: l'interfaccia comune, con le operazioni e (nella forma trasparente) la gestione dei figli.
\textbf{Leaf}: un oggetto senza figli. \textbf{Composite}: memorizza i figli e implementa le operazioni chiamandole su
di loro. \textbf{Client}: lavora solo con il Component, senza sapere se ha in mano una foglia o un composto.
La relazione decisiva è l'aggregazione \textbf{dal Composite al Component}: è lei che rende l'albero ricorsivo.
\end{struttura}

\begin{programma}{composite-esempio}{L'esempio del corso: il file system}
\begin{lstlisting}[style=javacs, name=EsempioComposite.java]
public class EsempioComposite {
    public static void main(String[] args) {
        Directory one = new Directory("dir111"),
                  two = new Directory("dir222"),
                  thr = new Directory("dir333");
        one.add(new File("a"));
        one.add(two);
        one.add(new File("b"));
        two.add(new File("c"));
        two.add(new File("d"));
        two.add(thr);
        thr.add(new File("e"));
        one.ls("");           // una sola chiamata
    }
}
interface AbstractFile {            // Component
    void ls(String rientro);
}
class File implements AbstractFile { // Leaf
    private final String nome;
    File(String nome) { this.nome = nome; }
    public void ls(String rientro) {
        System.out.println(rientro + nome);
    }
}
// Composite
class Directory implements AbstractFile {
    private final String nome;
    private final List<AbstractFile> figli =
            new ArrayList<>();
    Directory(String nome) { this.nome = nome; }
    void add(AbstractFile f) { figli.add(f); }
    public void ls(String rientro) {
        System.out.println(rientro + nome);
        for (AbstractFile f : figli)
            f.ls(rientro + "   ");   // ricorsione
    }
}
\end{lstlisting}
\tcblower
\begin{tikzpicture}
  \umlclass[2.6cm]{a}{(2.0,0)}{\stereotipo{interface}{AbstractFile}}{}{+ ls(rientro: String)}
  \umlclass[2.9cm]{f}{(-0.2,-2.6)}{File}{- nome: String}{+ ls(rientro: String)}
  \umlclass[3.4cm]{d}{(3.6,-2.6)}{Directory}{- nome: String}{\textasciitilde\ add(f: AbstractFile)\\+ ls(rientro: String)}
  \draw[realizza] (f.north) |- ++(0,0.5) -| ([xshift=-0.4cm]a.south);
  \draw[realizza] (d.north) |- ++(0,0.35) -| ([xshift=0.4cm]a.south);
  \draw[aggr] (d.east) -- ++(0.45,0) |- (a.east);
  \node[molt, anchor=south west] at ($(a.east)+(0.1,0.05)$) {* \{ordered\}};
  \node[ruolo, anchor=north west] at ($(a.east)+(0.1,-0.05)$) {figli};
\end{tikzpicture}

\medskip
\begin{lstlisting}[style=shellcs, name=Terminale]
$ javac EsempioComposite.java
$ java EsempioComposite
dir111
   a
   dir222
      c
      d
      dir333
         e
   b
\end{lstlisting}
\regola{\code{Directory} non sa se un figlio è un file o un'altra directory: chiama \code{ls} e ci pensa il
\textbf{polimorfismo}. Il \code{main} chiama \code{ls} una volta, sulla radice, e la ricorsione visita l'albero. Nel
codice del corso il rientro è una \code{StringBuffer} statica che ogni \code{ls} allunga e poi accorcia, e
\code{add} riceve un \code{Object}; qui il rientro è un parametro e \code{add} accetta solo \code{AbstractFile}.}
\end{programma}
